\section{An exact repair of the logical score optimization}
\label{sec:repair}

The output magnitudes are free optimization variables even when all
labels are fixed. For the scalar weights below, the cheapest competitors
use a negative overshoot on one fixed logical coset. We first identify
that optimizer in a finite-dimensional weighted Fourier problem, then
establish the scalar estimates used by the neural construction.

For a score $g$ on $G=\mathbb F_2^m$, write
$B_w(g)=\sum_s w_s|\widehat g(s)|$ with normalized Walsh coefficients.
The sharp atomic theorem identifies this weighted norm with the
variation cost on the retained physical code by signed input-flip
averaging of its equality atoms.

\subsection{A reduction to nonnegative Fourier coefficients}

Let $G=\F_2^m$, $M=|G|$, and let $w_s\ge w_0>0$. Continue to write
$B_w(g)=\sum_s w_s|\widehat g(s)|$ and
$W(x)=M^{-1}\sum_s w_s(-1)^{s\cdot x}$.

\begin{lemma}[The point-positive margin problem]\label{lem:positive-coeff}
Suppose $M\ge4$. Among scores satisfying $g(0)\ge1/2$ and
$g(x)\le-1/2$ for every $x\ne0$, there is a minimum-cost score with
$g(0)=1/2$, $\widehat g(0)<0$, and $\widehat g(s)\ge0$ for $s\ne0$.
If $w_s>w_0$ for every $s\ne0$, every optimizer has these properties.
\end{lemma}
\begin{proof}
First impose $c_s=\widehat g(s)\ge0$ for $s\ne0$ and
$c_0=1/2-\sum_{s\ne0}c_s$. The margin constraints become
\[
 \sum_{s:s\cdot x=1}c_s\ge\frac12\quad(x\ne0),
\]
and their sum implies $\sum_{s\ne0}c_s\ge1-1/M$, hence $c_0<0$.
The resulting objective is
$-w_0/2+\sum_{s\ne0}(w_0+w_s)c_s$.
This is a feasible, attained finite linear program. Its dual has
variables $a_x\ge0$ and constraints
\[
 \sum_{x:s\cdot x=1}a_x\le w_0+w_s,
\]
with objective $\frac12\sum_{x\ne0}a_x-w_0/2$.
For an optimal dual, define a function $\lambda$ by
\[
 \lambda(0)=M\left(\frac12\sum_{x\ne0}a_x-w_0\right),\qquad
 \lambda(x)=-Ma_x/2\quad(x\ne0).
\]
Then $\widehat\lambda(0)=-w_0$ and
\[
 \widehat\lambda(s)=-w_0+\sum_{x:s\cdot x=1}a_x\in[-w_0,w_s].
\]
Consequently $|\widehat\lambda(s)|\le w_s$. The restricted optimum
is at least $(3/2-2/M)w_0>w_0/2$, so $\lambda(0)>0$.
Thus $\lambda$ is also a dual certificate for the unrestricted
margin problem and has the same value. Equality forces $g(0)=1/2$
and $c_0\le0$. Under strict $w_s>w_0$, it also excludes negative
$c_s$ for $s\ne0$; the summed constraints then give $c_0<0$.
\end{proof}

If the prescribed positive set is a subgroup $J\le G$, averaging over
translations by $J$ deletes all Fourier coefficients outside $J^\perp$.
It strictly decreases the norm whenever any such coefficient is nonzero.
Thus every optimizer is $J$-invariant, and the problem reduces to the
point-positive problem on $G/J$ with weights restricted to $J^\perp$.

\subsection{A dominant Walsh coefficient selects a unique overshoot}

\begin{theorem}[Exact subgroup-positive optimizer]\label{thm:dominant}
Let $J\le G$, let $u\notin J$, and set $H=\langle J,u\rangle$,
$j=|J|$, and $L=M/(2j)\ge2$. Assume $w_s>w_0>0$ for $s\ne0$.
On $G/J$ put
\[
 A_x=\sum_{a\in J}W(x+a),\qquad t=\frac{2jw_0}{M},\qquad
 h_x=A_x-t,\qquad h=h_u>0.
\]
For every nonzero coset $c\in G/H$, choose $x\in c$ and put
\[
 \kappa_c=\max(0,A_x,A_{x+u}),\qquad
 r_c=\frac{\kappa_c}{h},\qquad R_* =\sum_{c\ne0}r_c.
\]
These definitions do not depend on the representatives. If
\begin{equation}\label{eq:dominant-conditions}
 R_*<1,\qquad w_0>(1+R_*)(h+t),
\end{equation}
then the unique minimum of $B_w(g)$ over the full-margin sign cell
positive exactly on $J$ is
\begin{equation}\label{eq:overshoot-score}
 g_J(x)=-\tfrac12+\one_{x\in J}-\one_{x\in u+J}.
\end{equation}
Its cost is
\begin{equation}\label{eq:overshoot-cost}
 C_J=\frac{w_0}{2}+A_0-A_u.
\end{equation}
\end{theorem}
\begin{proof}
Work on $G/J$. By Lemma~\ref{lem:positive-coeff}, every optimizer
has value $1/2$ at the origin and nonnegative nonconstant Fourier
coefficients. It can therefore be written uniquely as
\[
 g=-\tfrac12+\one_0-\sum_{x\ne0}\nu_x\one_x,\qquad \nu_x\ge0.
\]
Nonnegativity of its nonconstant coefficients is exactly
\begin{equation}\label{eq:overshoot-primal}
 \sum_{x\ne0}\nu_x(-1)^{s\cdot x}\le1
 \quad(s\in J^\perp,\ s\ne0).
\end{equation}
Substitution in the objective gives the equivalent maximization
\[
 C_J=\frac{w_0}{2}+A_0-t
       -\max_{\nu}\sum_{x\ne0}\nu_xh_x.
\]
We certify that the unique maximizer is $\nu_u=1$.

Define a probability law on $H^\perp$ by
\[
 P(s)=\frac1L\left[1+\sum_{c\ne0}r_c(-1)^{s\cdot x_c}\right].
\]
It is strictly positive because $R_*<1$, and character orthogonality
shows that it sums to one. Write $p=P(0)=(1+R_*)/L$ and condition on
$s\ne0$ to obtain $Q$. The second condition in
\eqref{eq:dominant-conditions} implies $p<1$. For $x\notin H$,
\[
 \E_Q(-1)^{s\cdot x}
 =\frac{r_{[x]}-p}{1-p}
 \ge\frac{A_x}{h}-\frac p{1-p}
 >\frac{A_x-t}{h}.
\]
The strict last inequality is equivalent to
$w_0>(1+R_*)(h+t)$, using $t=w_0/L$.
For $x=u$ modulo $J$, the character is identically one and $h_x/h=1$.
Averaging \eqref{eq:overshoot-primal} against $Q$ now gives
\[
 \sum_{x\ne0}\nu_xh_x
 \le h\E_Q\sum_{x\ne0}\nu_x(-1)^{s\cdot x}\le h.
\]
Equality requires $\nu_x=0$ outside $u+J$; the remaining constraint
forces $\nu_u\le1$. The feasible choice $\nu_u=1$ attains equality,
proving uniqueness and both displayed formulas.
\end{proof}

\subsection{Walsh signs from the inverse Laplace kernel}

The scalar construction only needs a positive inverse Laplace kernel
near zero. No derivative estimate whose order grows with the logical
dimension is required. Lemma~\ref{lem:laplace-kernel}, applied to
Theorem~\ref{thm:curvature}, gives
\[
 w(n)=An+B+\int_0^\infty e^{-nt}\rho(t)\,dt,
\]
where, for fixed positive constants, $c_0\le\rho(t)\le C_0$ on
$(0,t_0]$ and $|\rho(t)|\le C_1e^{c_1t}$ everywhere.

\begin{lemma}[Uniform high-order Walsh signs]\label{lem:uniform-walsh}
Consider an auxiliary increment cube with $q$ coordinates, increments
$\delta_i\ge2$, total increment $D_\delta=\sum_i\delta_i$, and
$n_s=n_0+\sum_i\delta_i s_i$. Along a sequence with $n_0\to\infty$,
assume
\[
 D_\delta=o(n_0),\qquad
 q\log(n_0+D_\delta)=o(n_0).
\]
For every support $T$ of size $2\le k\le q$, uniformly in $T$,
\begin{equation}\label{eq:uniform-walsh}
 \frac{c\prod_{i\in T}\delta_i}
      {2^k(n_0+D_\delta)^{k+1}}-E_0
 \le W(T)\le
 \frac{C k!\prod_{i\in T}\delta_i}{2^k n_0^{k+1}}+E_0,
 \qquad E_0=O(e^{-c'n_0}).
\end{equation}
In particular all these coefficients are strictly positive for all
sufficiently large construction parameters. Singleton coefficients
are negative.
\end{lemma}
\begin{proof}
The affine part $An+B$ has zero coefficient at $|T|\ge2$.
Factoring the Laplace integrand gives
\[
 W(T)=\int_0^\infty e^{-n_0t}\rho(t)
 \prod_{i\in T}\frac{1-e^{-\delta_it}}2
 \prod_{i\notin T}\frac{1+e^{-\delta_it}}2\,dt.
\]
All factors except $\rho$ lie in $[0,1]$. For $a,t>0$,
\[
 (at/2)e^{-at}\le(1-e^{-at})/2\le at/2,
 \qquad e^{-at}\le(1+e^{-at})/2\le1.
\]
Put $N=n_0+D_\delta$. On $[1/N,2/N]\subset(0,t_0]$, these
inequalities give a positive contribution at least
$c_0e^{-2}\prod_{i\in T}\delta_i/(2^kN^{k+1})$.
The potentially adverse tail on $t\ge t_0$ is bounded in absolute
value by $E_0=O(e^{-c'n_0})$. This proves the lower bound.
For the upper bound, use $\rho\le C_0$ near zero and extend the
nonnegative majorant $t^ke^{-n_0t}$ to the full half-line. Its gamma
integral is $k!/n_0^{k+1}$, giving the upper bound.
Since $\delta_i\ge2$, the positive lower contribution is at least
$cN^{-(q+1)}$, whose logarithm is $-o(n_0)$ by hypothesis.
It therefore dominates $E_0$ uniformly, proving positivity.
Finally $w'(a)=A+O(a^{-2})>0$ for large $a$, and the first-difference
formula proves the singleton assertion.
\end{proof}

\begin{lemma}[A dominant pair of increments]\label{lem:dominant-increments}
Under Lemma~\ref{lem:uniform-walsh}, let two distinguished increments
belong to $[H,2H]$ and put $S=\sum_{j\ge3}\delta_j$.
If $S=o(H)$ and $H=o(n_0)$, then, for $u=e_1+e_2$,
\begin{equation}\label{eq:dominance-sum}
 \frac{\sum_{|T|\ge2,\ T\ne u}W(T)}{W(u)}
 =O(S/H+H/n_0).
\end{equation}
\end{lemma}
\begin{proof}
The lower bound gives $W(u)\ge cH^2/n_0^3$ for large parameters.
Other degree-two increment products sum to $O(HS+S^2)$, so their
upper bounds divided by $W(u)$ give $O(S/H)$, apart from tails.
Write $D_\delta=O(H)$. For $k\ge3$ use
$k!\sum_{|T|=k}\prod_{i\in T}\delta_i\le D_\delta^k$ and sum
$\sum_{k\ge3}D_\delta^k/(2^k n_0^{k+1})=O(D_\delta^3/n_0^4)$.
Dividing by $W(u)$ gives $O(H/n_0)$.
The total tail is at most $2^qE_0$, exponentially negligible relative
to $W(u)$ because $q\log(n_0+D_\delta)=o(n_0)$.
\end{proof}

\paragraph{Splitting large increments among repeated columns.}
The physical atom theorem requires each individual rail increment to
be at most its baseline $R$. A large effective increment $H$ can
nevertheless be realized using repeated copies of one logical column.
For even integers $H,R$, write $H$ as a sum of
$\lceil H/R\rceil$ positive even integers, each at most $R$.
Assign those increments to that many identical columns. Every retained
transversal chooses all their long rails together, so their scalar
degree contribution is exactly $H(a\cdot s)$.
Thus before the scalar Walsh analysis these copies may be collapsed
into one effective increment, while the physical atom theorem keeps
them separate. Repeated column directions do not change the Cayley
independent-set condition used for observability.

The construction uses $R=\Theta(\log^2m)$ and
$H=\Theta(m\log m)$. It has $O(m)$ ordinary physical columns and
$O(H/R)=O(m/\log m)$ additional copies in each distinguished group.
Hence the total pair count is $p=\Theta(m)$ and
$n_0=pR+1=\Theta(m\log^2m)$. On the auxiliary cube there are
$q=O(m)$ effective columns, $S=O(m)$, and
$D_\delta=O(m\log m)=o(n_0)$; all hypotheses above hold.
After restriction to the logical column code, higher-order coefficients
are pushed forward as sums. Negative singleton contributions occur only
at individual columns. Neither the distinguished pair $u$ nor any live
candidate is such a column, so the required dominant coefficient and
sign conditions are preserved.
